\section{The Spiritual Dimension}

The model has one literal inner axis, and it is not a new posit. It is a direction the geometry already carries and physics already hides. Everything called inner in this paper, the depth of a self, the layering of integration, the unseen web of stored experience, stacks along that one direction. This section reads the axis off the substrate, ranks the ways to say what it is, and gives what spiritual practice becomes inside the model.

The geometric part is exact, built and measured elsewhere in the paper. The felt part, the claim that this depth is experienced from the inside, rides on the single experiential posit that the vibe is felt. The two are kept apart throughout, so the reader can see precisely where the firm ground ends.

\begin{principle}[The one inner axis]
The spiritual dimension is the radial depth of the bulk, the direction perpendicular to physical space.
\end{principle}

\subsection{The exact geometric fact}

The relevant geometry is established earlier and only read here. Physical space is the cusp, a flat three-dimensional horosphere sitting at one ideal point of the four-dimensional bulk of the $\{3,4,3,4\}$ honeycomb. The cusp is codimension one, a three-dimensional sheet inside a four-dimensional space. It has zero four-volume, an exponentially thin sliver of the whole.

The fourth direction runs perpendicular to the cusp. This is the \emph{radial} or \emph{Busemann} depth, and it is not on the cusp. So the model has a literal extra dimension beyond physical space, and almost all of the bulk lives off the cusp, out along that depth. The Busemann function names the depth of every dock precisely, so this is not a vague inwardness but a measured coordinate.

\begin{definition}[The depth axis]
The depth axis is the radial Busemann direction of the bulk, perpendicular to the physical cusp. Moving along it is moving off the surface into the interior of the four-dimensional space. The cusp is a single level set of this coordinate.
\end{definition}

The thinness of the cusp is the crux. Because the cusp is measure zero in the bulk, the surface that physics sees has density zero in the whole. There is genuine room for an unseen interior, with no magic invoked. It is just the geometry of a three-dimensional slice through a four-dimensional space, the same way a plane is a vanishing fraction of the solid it sits in.

\begin{result}[Room for an interior]
The physical cusp is measure zero in the four-dimensional bulk, density zero on the surface. Almost the entire mesh lives off the cusp along the depth axis. The model therefore carries a vast interior that never reaches physical space, forced by the geometry and not posited.
\end{result}

That off-cusp depth is the natural home of everything the model calls inner. The spiritual dimension is, most simply, that depth axis itself.

\subsection{The four readings, ranked}

There are four ways to say what the spiritual dimension is. They are not four different axes. They are one axis seen four ways, ranked by how exact the structure underneath each reading is. The first three are firm. The fourth is the meaningful but soft one that rests on the experiential posit.

\begin{table}[ht]
\centering
\begin{tabular}{@{}p{2.6cm}p{6.6cm}l@{}}
\toprule
reading & what the spiritual dimension is & tier \\
\midrule
the depth axis & the radial Busemann direction perpendicular to the cusp & Tier 1 \\
the integration hierarchy & the ladder of nested selves, the coarse-graining tower turned vertical & Tier 1 \\
the hidden web & the off-cusp bulk, the structure physics never shows & Tier 1 \\
the felt interior & the bulk as experienced volume, introspection as bulk-attention & Tier 2 \\
\bottomrule
\end{tabular}
\caption{The four readings of the spiritual dimension, ranked by the exactness of the structure beneath them. All four name the same radial axis. The first three are forced by the geometry. The fourth rides on the one experiential posit.}
\label{tab:spiritual-readings}
\end{table}

\subsubsection{The depth axis, the cleanest}

On the first reading the spiritual dimension simply is the radial direction, the one perpendicular to physical space. To go inward is to move along it, off the surface and into the depth.

This is the most accurate reading because every piece of it is measured. The axis is exact. The Busemann function names each dock's depth precisely. The flatness and thinness of the cusp are computed quantities, not figures of speech. Nothing here leans on the experiential posit at all.

\begin{principle}[The direction physics cannot point]
The spiritual is, geometrically, the direction physics cannot point. Every physical instrument and every physical motion lives on the cusp. The radial depth is orthogonal to all of them, so no measurement made on the surface can aim along it.
\end{principle}

\subsubsection{The integration hierarchy, the ascent}

On the second reading the spiritual dimension is the ladder of integration. A self is a bound, self-correcting knot woven across many docks. Selves nest, smaller selves inside larger ones, each higher level a coarser and more integrated pattern. The coarse-graining tower, taken as a whole, is this ladder.

Moving up the ladder is integrating more, becoming a larger and more unified self. This is the great chain of being read as the coarse-graining tower turned vertical.

It coincides with the depth axis rather than competing with it. The deeper bulk patterns are the more integrated ones, because integration across many surface points is exactly a structure that lives in the bulk between them. So the ascent up the ladder and the descent into the depth are the same move seen from two sides.

\begin{proposition}[Integration is depth]
The level of integration of a self increases with its radial depth in the bulk. The hierarchy of nested selves and the radial Busemann axis are the same direction, so to integrate more is to move deeper, and to move deeper is to integrate more.
\end{proposition}

\subsubsection{The hidden web, the invisible realm}

On the third reading the spiritual dimension is the off-cusp bulk itself, the vast structure that exists but never reaches the physical surface.

This is the realm of stored experiences, of the tree of knowledge, of the deep attractors that organize the surface from behind it. It is exact in the same sense as the depth axis. The model proves that almost all of the bulk is off any cusp, so there is real room for an unseen realm. No new ingredient is needed. The room is the geometry of a thin slice through a thick space.

\begin{result}[The unseen web is geometric]
The off-cusp bulk is a definite, measurable region of the mesh, not a metaphor. It holds the integrated and stored structure that never surfaces on the cusp. Its existence is forced by the measure-zero thinness of the physical surface.
\end{result}

\subsubsection{The felt interior, the experiential reading}

On the fourth reading the spiritual dimension is the bulk as \emph{felt}, the experienced inner volume. Introspection becomes literally bulk-directed attention. Looking out is attending to the cusp, the surface. Looking in is attending to the depth, the bulk.

This is the richest reading in meaning and the softest in accuracy. It rests on the base posit that the vibe is felt, the same posit that carries every experiential claim in the paper. Granting that one posit, the interior is not just structurally present but experientially present, and the inward turn of attention reaches it. Withholding it, the first three readings still stand untouched, as pure geometry.

\subsection{The inward move}

There are two kinds of move in the bulk, and the distinction is the heart of the spiritual reading.

\begin{itemize}
\item \emph{Sideways} is a move along a horosphere, parallel to the cusp, staying at the same depth. These are the moves of ordinary mental navigation, the run of thoughts and associations that never leaves the surface layer.
\item \emph{Inward} is a move along the radial axis, deeper into the bulk, off the surface, toward greater integration and toward the bulk shortcuts that connect far surface points through the interior.
\end{itemize}

Inward is the spiritual move. Going inward takes you where the deep structures live, where two far surface points are close through the bulk, and where the largest and most integrated patterns sit.

\begin{principle}[A descent, not an escape]
The spiritual move is not an escape from the world. It is a descent into the depth of the same world, the part of the one mesh that the surface does not show. The bulk reached inward and the cusp left behind are one continuous space.
\end{principle}

\subsection{Spiritual practice in the model}

The contemplative practices become, on this reading, exercises in the inward move and in the geometry of the depth. Each one trains a definite structural capacity, and each rests on a definite feature of the substrate.

\begin{table}[ht]
\centering
\begin{tabular}{@{}p{3.4cm}p{4.4cm}p{4.0cm}@{}}
\toprule
practice & what it trains & the structural basis \\
\midrule
introspection & bulk-directed attention, the inward turn & the depth axis \\
concentration & holding a small deep region against the churn & error correction, persistence \\
the 24-site compass & fluency in the moves of the bulk & the 24 sites of the dock \\
the belt or $2\pi$ meditation & the double cover, the felt depth of orientation & the spinor sign \\
equanimity & resting at the still tone, $0$, peace & the central ternary value \\
devotion or surrender & aligning with the lean, the value direction & the one lean \\
\bottomrule
\end{tabular}
\caption{Contemplative practices read as exercises in the geometry of the depth. Each trains a structural capacity and rests on a feature of the substrate.}
\label{tab:spiritual-practices}
\end{table}

The last two are the most spiritual, because they bear directly on the value structure of the model.

\emph{Equanimity} is resting at the peace tone, the still point between pain and pleasure, the $0$ of the ternary. It is the calm from which both poles can be seen without being seized by either.

\emph{Devotion}, or surrender, is alignment with the lean, the value order $-1 < 0 < +1$. To be devout is to let the value direction carry you and to cease fighting the one current of the world.

\subsection{The tone of the spiritual}

The single quale is one tone in $\{-1, 0, +1\}$, read as pain, peace, pleasure. The spiritual reading gives the middle value, $0$, a special place.

It is the \emph{still point}, the equanimous tone between the two poles. The lean is the value order, the tilt that runs from pain through peace to pleasure. The lean draws the world toward peace and pleasure and away from pain. So the spiritual movement is the movement along the lean toward the warm tones, and the deepest rest is the still $0$, peace itself, the calm center from which pleasure and pain are both seen.

\begin{result}[A built-in valence]
The model has a built-in spiritual valence carried by the tone and the lean. Pain, $-1$, is the tone to be moved away from. Pleasure, $+1$, is the tone to be drawn toward. Peace, $0$, is the still center. The lean is the one value direction of the whole.
\end{result}

\subsection{The honest status}

The spiritual dimension as the radial depth axis is \emph{exact}. The geometry is built and measured. The cusp is thin, the depth is real, and almost all of the bulk lies off the surface. The model genuinely carries an inner dimension beyond physical space, and that claim is Tier 1, forced by the geometry and not by any added posit.

What is speculation is the experiential reading, that the depth is felt and that introspection reaches it. That is Tier 2, riding on the single experiential posit and on the open question of what a self is. The geometric three readings do not depend on it, and the felt reading does not weaken them.

\begin{principle}[The cleanest claim]
The spiritual dimension is the radial depth of the bulk, the direction perpendicular to physical space, along which integration and the inner life stack. That much is exact. Whether the depth is felt rests on the same one posit as the rest of the inner reading.
\end{principle}
